\begin{pastexam}[label={px:gk3-099-ds-03}]{高考三級資訊處理 99 年 第三題}
（申論；20 分）

\begin{itemize}[leftmargin=2.2em]
\item[㈠] 說明樹（tree）與二元樹（binary tree）有那三項主要的不同？（5 分）
\item[㈡] 已知某一樹其分支度（degree）為 1 的節點（node）有 5 個，分支度為 2 的節點有 4 個，分支度為 3 的
節點有 3 個，分支度為 4 的節點有 2 個，分支度為 5 的節點有 1 個，請問此樹一共有幾個節點？（5 分）
\item[㈢] 證明：於任意一個二元樹中，若$n_0$代表分支度為 0 的節點數目，$n_1$代表分支度為 1 的節點數目，
$n_2$代表分支度為 2 的節點數目，則$n_0 = n_2 + 1$。（10 分）
\end{itemize}
\end{pastexam}

\begin{pastsoln}
答案：㈠ 二元樹可以是空的、樹不行；二元樹每個節點最多兩個子節點、樹沒有限制；二元樹的子節點分左右、樹
的子節點不分。 ㈡ $36$ 個。 ㈢ 見下面的證明。

\textbf{㈠} 三項主要的不同（資結書第~7.1~節）：
\begin{enumerate}
\item \textbf{能不能是空的}：樹至少有一個節點（根）；二元樹可以是空樹（沒有任何節點）。
\item \textbf{分支度的上限}：樹的節點可以有任意多個子節點；二元樹每個節點最多兩個。
\item \textbf{子節點有沒有順序}：樹的子樹之間不分次序；二元樹的子樹分成左子樹和右子樹，只有一個子節點時
也要分是左還是右。 例如「根 A 只有左子節點 B」和「根 A 只有右子節點 B」是兩棵不同的二元樹，當成樹看卻
是同一棵。
\end{enumerate}

\textbf{㈡} 設總節點數 $n$，邊數 $e$。 每個節點的分支度就是它的子節點數，每條邊恰好連到一個子節點，所以
邊數＝所有分支度的和：
\[
e = 1 \cdot 5 + 2 \cdot 4 + 3 \cdot 3 + 4 \cdot 2 + 5 \cdot 1 = 5 + 8 + 9 + 8 + 5 = 35.
\]
樹有 $n$ 個節點就有 $n - 1$ 條邊（除了根，每個節點恰好有一條連到父節點的邊），所以 $n = e + 1 = 36$。
其中分支度為 $0$ 的葉節點有 $36 - (5 + 4 + 3 + 2 + 1) = 21$ 個。

\textbf{㈢ 證明。} 設二元樹有 $n$ 個節點。 若 $n = 0$，題目的式子不成立（$0 \ne 0 + 1$），所以這裡的「任意
一個二元樹」指非空的二元樹，$n \ge 1$。

\textbf{數節點。} 每個節點的分支度只能是 $0$、$1$、$2$，所以
\[
n = n_0 + n_1 + n_2. \tag{1}
\]
\textbf{數邊。} 除了根以外，每個節點恰好有一條邊連到它的父節點，所以邊數是 $n - 1$。 另一方面，每條邊都
是從某個節點連到它的一個子節點，分支度 $1$ 的節點貢獻 $1$ 條、分支度 $2$ 的貢獻 $2$ 條，所以邊數也等於
$n_1 + 2n_2$：
\[
n - 1 = n_1 + 2n_2. \tag{2}
\]
把 (1) 代入 (2)：$n_0 + n_1 + n_2 - 1 = n_1 + 2n_2$，消去 $n_1$，得到 $n_0 = n_2 + 1$。 \hfill$\blacksquare$

這就是資結書定理~7.2。 注意結果和 $n_1$ 無關：在一條鏈上多插幾個只有一個子節點的節點，葉子數和分支度 $2$
的節點數都不變。 ㈡ 的一般樹也可以用同樣的方法算：$n - 1 = \sum_k k \cdot n_k$。
\end{pastsoln}
